% PONTE ENTRE O MODELO E O LIVRO
% ==============================
% Público: autores usam \aberturaartigo dentro de artigos/modelo_artigo/.
% Os campos vêm de metadados.tex. Não edite esta macro para um artigo isolado;
% isso causaria diferenças de layout entre capítulos.

\newcommand{\aberturaartigo}{%
  \capitulocapa
    {\ArtigoTitulo}
    {\ArtigoSubtitulo}
    {\ArtigoAutores\\ \ArtigoAfiliacoes\\ \ArtigoContatos}
    {\textbf{Resumo.} \ArtigoResumo}
    {\ArtigoAutores}%
  % \capitulocapa encerra a página de abertura, que é a única com o
  % mini-cabeçalho. As páginas seguintes usam o cabeçalho corrente do artigo.
  \pagestyle{fancy}%
  \noindent\textbf{Palavras-chave:} \ArtigoPalavrasChave\par\vspace{0.8cm}
  \begin{center}
    {\fontsize{14}{16}\selectfont\color{sommagreen}\textbf{\ArtigoTituloIngles}}
  \end{center}
  \noindent\textbf{Abstract.} \ArtigoAbstract\par\vspace{0.4cm}
  \noindent\textbf{Keywords:} \ArtigoKeywords\par\vspace{0.8cm}
}

% O modelo deixa as referências por último e evita que autores repitam a mesma
% configuração de biblatex em cada artigo.
\newcommand{\referenciasartigo}{%
  \section*{Referências}
  \printbibliography[heading=nossabibliografia]
}

% Figura no padrão já usado pelo template_tcc_monografia: ambiente figure[H],
% legenda acima, label para referência, imagem centralizada e fonte abaixo.
% Argumentos: [largura], rótulo, título, arquivo e fonte.
% Exemplo: \inserirfigura[0.45\textwidth]{fig:prot}{Protótipo}{imagens/a.png}{Elaboração própria (2026).}
\newcommand{\inserirfigura}[5][0.65\textwidth]{%
  \begin{figure}[H]
    \centering
    \caption{\label{#2}#3}
    \includegraphics[width=#1]{#4}
    \par\vspace{0.15cm}
    {\fontsize{11}{13}\selectfont Fonte: #5\par}
  \end{figure}
}
